% \chapter{Self-Organising Principles of Intelligence: Self-Organising Digital Circuits}
\chapter{Self-Organising Digital Circuits}
\label{chap4}

\begin{abstract}

\subsection*{Narrative Preamble}

If \Cref{chap3} explored the computational potential of local rules, this chapter explores their adaptive potential. We move closer to the biological reality of a system that must not only compute, but also maintain itself and adapt to changes.
Here, we utilize the principles of self-organization as a "meta-learning" mechanism. Instead of the NCA being the computer, it becomes the optimizer. We investigate a system where local rules navigate and configure a functional digital circuit, effectively replacing global backpropagation with a decentralized, graph-based optimization process. This architecture has implications for true scalable intelligence: a system where the "learning rule" is itself a local, robust policy. By demonstrating how such a system can recover from physical damage and reconfigure itself without global supervision, we provide a concrete instantiation of the "adaptive resilience" that characterizes biological intelligence, pointing the way toward systems that truly know how to learn and adapt.

\subsection*{Publication Context}

This chapter presents ongoing, unpublished collaborative work with Marcello Barylli from the group of Prof. Sebastian Risi at the IT University of Copenhagen. The project originated from a concept proposed by Alexander Mordvintsev at the 2025 Capocaccia Neuromorphic Intelligence Workshop.

\subsection*{Contributions}

Alexander Mordvintsev conceived the initial project idea and developed the inital software: the implementation of differentiable Boolean circuits and the graphical interface for observing standard backpropagation behaviour. The author and Marcello Barylli formalized the architectural shift to use a Transformer as a meta-optimizer. The author led the development and optimization of the NCA framework, implementing graph-based models, the meta-learning infrastructure and the random pool mechanism, and conducted the bulk of the "growth" wiring experiments. Marcello conducted the primary experiments regarding functional maintenance, and the analysis on parameter space trajectories.

\end{abstract}

\section{Abstract}
% Abstract length should not exceed 250 words
Fault tolerance in classical computing has traditionally relied on static strategies like hardware redundancy and error-correcting codes. Biological systems, in contrast, exhibit \textit{adaptive plasticity}, maintaining function through dynamic re-organisation around damage. Inspired by this principle, we introduce a framework for self-organising digital circuits that replaces rigid redundancy with learned, autonomous self-repair. We extend the \acf{NCA} paradigm from pattern formation on grids to functional logic generation on arbitrary graphs. Our architecture employs a topology-masked Transformer as a decentralized meta-optimizer that learns to configure the \acfp{LUT} of \acfp{LGN}. Unlike standard NCA which regenerate a fixed target state, our model learns a functional local policy that navigates the Boolean search space to satisfy a computational task. We demonstrate that this policy can self-assemble functional circuits from scratch and, crucially, rapidly \textit{re-route} logic around permanent, previously unseen hardware faults. This work bridges the principles of biological self-organisation with the practical domain of digital hardware, offering a blueprint for resilient systems capable of graceful degradation and real-time adaptation.

\footnote{Data/Code + \textbf{Circuit Visualization} available \href{https://github.com/GabrielBena/boolean_nca_cc/tree/gabi}{here}}


\section{Introduction}

% Motivation: Engineering vs Biological Fault Tolerance
\subsection*{Motivation: From Prescriptive Redundancy to Adaptive Plasticity}
Robustness is a central tenet of modern informatics. Decades of research have yielded systems capable of impressive fault tolerance, utilizing mechanisms such as Error Correcting Codes (ECC), modular redundancy, and sophisticated fallback protocols to ensure graceful degradation \citep{neumann_probabilistic_1956, schuster_robustness_2008}. However, these engineered successes typically rely on \textit{anticipation}: resources are pre-allocated and failure modes are modelled in advance. When a system encounters damage that exceeds its redundancy budget or defies its failure model, it might fail brittly.
Biological intelligence offers a complementary paradigm: \textit{adaptive plasticity}. Natural automata, such as the mammalian cortex, do not rely solely on static backups. Instead, they exhibit the ability to dynamically re-purpose surviving components to compensate for injury, a phenomenon known as cortical remapping \citep{nudo_recovery_2013}. This form of resilience is not pre-programmed but emergent; the system "learns" its way around the damage.
Inspired by these biological principles, and by early von Neumann's assumption that systems should "operate across errors" \citep{von_neumann_theory_1966}, we seek to endow digital hardware with similar capabilities. This work introduces a framework for self-organising digital circuits that complements traditional redundancy with dynamic, learned self-repair. 
In this sense, it is a new take on the idea of self-repairing digital hardware, which has already been considered as a robustness mechanism, especially in FPGAs \citep{carter_user_1986, schuster_robustness_2008, gokhale_dynamic_2004}

% Fault tolerance in classical computing has traditionally relied on rigid architectural strategies, such as hardware redundancy and error-correcting codes. While effective, these methods depend on static pre-allocation of resources—essentially anticipating failure before it occurs. 
% Biological intelligence, in contrast, exhibits mechanisms that allow for \textit{adaptive resilience}. Natural automata, such as the brain, do not rely solely on backup components; rather, they maintain function through dynamic re-organisation around damaged areas, a phenomenon well-documented in neural recovery studies \citep{nudo_recovery_2013}.
% Inspired by von Neumann's early postulates on systems that "operate across errors" \citep{von_neumann_theory_1966}, we seek to bridge this gap between biological plasticity and digital hardware. This chapter introduces a framework for self-organising digital circuits that replace static redundancy with dynamic, learned self-repair.

\subsection*{The Evolution of Automata: From Fixed Rules to Learned Dynamics}
% CAs intro: Complexity and Turing Completeness
The study of \acf{CA} has long served as a fertile ground for investigating the essence of computation and complexity. Fundamentally, a Cellular Automaton is a discrete computational system where a grid of cells, each holding a specific state, evolves in synchronized steps. The evolution of each cell follows the same local rule that dictates how every cell updates based on its immediate neighbours on the grid. From simple local rules, CAs can generate remarkably intricate patterns and have been proven to be Turing complete \citep{cook}. 
% CAs limits: Hard to find rules
However, CAs suffer from a significant control problem: while it is easy to observe complexity arising from pre-determined rules \citep{wolfram}, explicitly engineering specific local rules to achieve a desired global outcome is notoriously difficult.
Attempts to extend these systems into the continuous domain, such as Lenia  or SmoothLife \citep{lenia, smoothlife}, introduced life-like fluidity but further complicated the stability and control of the system's outputs.

% NCA intro: The Paradigm Shift
Recently, \acf{NCA} have emerged as a paradigm shift that bridges the gap between self-organising systems and differentiable learning. In an NCA, the update rule is no longer hard-coded; it is a parameterised neural network \citep{mordvintsev_growing_2020}.
Each cell in the system utilizes a "perception" network to gather information from its local neighbourhood (typically on a grid), and an "update" network to compute an incremental change to its state. 
Because these operations are differentiable, the dynamics of the system can be "steered" via gradient descent towards specific goals. 
Traditionally, NCA have been utilized as models of morphogenesis: for example growing target images by treating the final cell states as RGB values. In these architectures, cells also maintain a vector of "hidden states," allowing the local rule to encode arbitrary information required to coordinate global structure. This distinction already let us differentiate between part of the cells' states that are "functional" (such as the RGB channels) and other that are "latents".

\subsection*{Beyond the Grid: Growing Function over Form}
While foundational, standard NCA remain topologically constrained to rigid grids and functionally constrained to pattern formation. To overcome these limitations, recent research has sought to extend the NCA paradigm into more flexible and general models of computation.
% 1. The Topological Shift: From Grid to Graph
First, the topological constraint: by reinterpreting the NCA update step as a message-passing operation on a graph, it is possible to generalise the concept of "neighbourhood" from spatial proximity to topological connectivity \citep{grattarola_learning_2021}. In this context, Graph Neural Networks (GNNs) provide the natural machinery for self-organisation: nodes "perceive" their neighbours via messages sent along edges, aggregate this information, and update their internal state (in an incremental manner, to stay true to the core NCA dynamics).
% 2. The Functional Shift: From Data to Logic
Second, the functional constraint: traditionally, NCA states represent the target pattern itself (e.g., RGB values), but this can also be extended. We've seen in \Cref{chap4} that we can move closer to general computation by interfacing continuous NCA with matrix tasks. While this demonstrated that NCA could learn a dynamic \textit{process} (e.g., matrix multiplication) rather than a static final state, it still conflated the CA's state with the computational data: effectively, the cells states were still holding the final result of the computation (the matrix values).
We propose here a more fundamental shift: viewing the cellular state not as the data to be processed, but as the \textit{functional substrate} that performs the subsequent processing.

% \subsection*{Differentiable Boolean Circuits}

% 3. The Implementation: Digital Boolean Circuits
To implement this, we utilise digital Boolean circuits as our substrate. Adopting the framework of Differentiable \acf{LGN} \citep{petersen_deep_2022}, we instantiate graphs where we treat nodes as gates defined by Lookup Tables (LUTs) and edges as functional wires.
Here, the "state" of a node is not a pixel colour or a matrix value, but the functional logic (the LUT configuration) it implements. By recursively applying a decentralised message-passing scheme, our NCA then acts as a meta-optimiser: it steers the substrate not toward a visual shape, but toward a configuration that is later-on able to implement a specific logic task (e.g., binary multiplication, addition, XOR etc).

\subsection*{The topology-masked Transformer}
% Meta-Learning Context: Masked Self-Attention and In-Context Learning
As we've just discussed, GNNs are well suited to update graphs residually, in a manner that extends the NCA paradigm to act on unconstrained topologies. Our architecture, however, employs a Transformer as the update rule. To see why this is natural, it is helpful to view self-attention through the lens of message passing. In a standard self-attention layer, each token updates its representation by computing a weighted sum over all other tokens, where the weights are determined by learned key-query interactions. This is, in essence, a message-passing operation on a \textit{complete graph}: every token sends a message to every other, and the receiver decides---via the attention weights---how much to listen to each sender. The crucial difference with classical GNNs is that the aggregation function is not fixed (e.g., sum or mean); instead, it is \textit{learned and input-dependent}, allowing the model to dynamically route information based on content.

Nonetheless, the graph-NCA paradigm requires strict locality: a node should only exchange information with its direct graph neighbours at each step. We enforce this by applying a \textit{topology mask} $M$ to the attention matrix, where $M_{ij}=1$ if and only if a wire connects gates $i$ and $j$ (or $i=j$), and $M_{ij}=0$ otherwise. Masked-out entries receive $-\infty$ before the softmax, ensuring zero attention weight. This single modification transforms the complete-graph communication pattern into message passing along the circuit's edges: each gate attends only to its wired neighbours and itself, with the key-query mechanism learning \textit{how} to weight and combine these neighbour messages at each step. We note that this is slightly distinct from ``Graph Transformer'' architectures \citep{dwivedi2021generalizationtransformernetworksgraphs}, which typically introduce graph-specific modifications such as structural encodings or edge features within the attention computation itself. Our approach is deliberately minimal: the topology enters \textit{only} through the attention mask, and all expressivity comes from the standard Transformer machinery. 
% This simplicity is a feature, not a limitation---it ensures that the learned update rule is a generic local policy, agnostic to any particular graph encoding scheme.

Furthermore, this design aligns with recent advances in meta-learning optimisers and general-purpose in-context learning \citep{kirsch_meta_2022, kirsch_general-purpose_2024}, which have demonstrated that Transformers can be meta-trained to internalise learning algorithms---effectively replacing explicit optimisers (like SGD) with a forward-pass policy.
Similarly in our system, the topology-masked Transformer effectively becomes the optimiser: a single, shared-weight attention block is applied recurrently for many steps, gaining representational depth through \textit{time} rather than through stacked layers. It learns a robust, local policy that navigates the landscape of Boolean configurations to find functional solutions.

\subsection*{Contributions and Key Findings}
% Synthesis and Hardware Implications
By applying these self-organising principles to digital hardware, we present a concrete instantiation of "adaptive resilience". Unlike standard backpropagation, which requires global knowledge of the circuit and differentiable components to repair damage \citep{sunada_blending_2025}, our system relies solely on local forward passes of the NCA.
This decoupling could allow for potential deployment on re-programmable hardware, such as FPGAs, where the "repairing agent" can operate efficiently on-chip.
We demonstrate that this meta-learned network can not only discover functions from scratch but also rapidly reconfigure circuits to recover from permanent, unseen physical damage, pointing the way toward autonomous systems that truly know how to learn and adapt.


\section{Methods}

\subsection{Differentiable Logic Gate Networks}
\label{sec:LGNs}
To enable gradient-based optimisation of digital logic, we adopt the \ac{LGN} framework \citep{petersen_deep_2022}.
A LGN is a Directed Acyclic Graph (DAG) where nodes represent logic gates and edges represent wires carrying bit signals. Unlike standard digital circuits where gates have fixed truth tables (e.g., AND, XOR), our gates are parameterised by learnable Lookup Tables (LUTs).

\begin{table}[h]
    \centering
    \caption{Standard Look-Up Tables (Truth Tables) for 2-Input Logic Gates}
    \label{cha5:tab:std_luts}
    \begin{tabular}{cc|cccccc}
        \toprule
        \multicolumn{2}{c}{\textbf{Inputs}} & \multicolumn{6}{c}{\textbf{Gate Output}} \\
        \cmidrule(lr){1-2} \cmidrule(lr){3-8}
        $A$ & $B$ & \textbf{AND} & \textbf{OR} & \textbf{XOR} & \textbf{NAND} & \textbf{NOR} & \textbf{L} \\
        \midrule
        0 & 0 & 0 & 0 & 0 & 1 & 1 & $l_{00}$ \\
        0 & 1 & 0 & 1 & 1 & 1 & 0 & $l_{01}$ \\
        1 & 0 & 0 & 1 & 1 & 1 & 0 & $l_{10}$ \\
        1 & 1 & 1 & 1 & 0 & 0 & 0 & $l_{11}$ \\
        \bottomrule
    \end{tabular}
\end{table}

\textbf{Continuous Relaxation of a gate:}
Standard Boolean operations are non-differentiable. We relax this by representing binary signals as continuous probabilities $x \in [0, 1]$.
For a gate with arity $N$ (number of inputs), the LUT is a tensor of $2^N$ learnable logits (such as the $L$ column of \cref{cha5:tab:std_luts}).
The "soft" output of a gate is computed via multilinear interpolation, effectively traversing the LUT as a probabilistic Binary Decision Diagram (BDD). For a gate with inputs $\mathbf{x} = [x_1, \dots, x_N]$ and LUT parameters $\mathbf{L}$, the forward pass recursively interpolates the table. For instance, with a 2-input gate, the first input $x_1$ interpolates between the two halves of the LUT to produce an intermediate tensor $L'$:
$$L^{\prime}=\left(1-x_1\right) \cdot\left[l_{00}, l_{01}\right]+ x_1 \cdot\left[l_{10}, l_{11}\right]$$
The second input $x_2$ does the same in $L^{\prime}$ and computes the final scalar output: 
$$y =\left(1-x_2\right)\cdot L^{\prime}[0]+x_2 \cdot L^{\prime}[1]$$

This operation is fully differentiable, allowing gradients from a loss on output $y$ to flow from back to the individual LUT entries. In practice, parameters are logits $\ell$, from which we get look-up table entries $L$ before functional execution. During training, we get $L = \sigma(\ell)$ with $\sigma$ being the sigmoid function. During inference however, the circuit can operate in "hard" mode by rounding LUT entries to $\{0, 1\}$ with $L = \Theta(\sigma(\ell))$, with $\Theta$ being a hard rounding function, thus recovering discrete Boolean logic.

\textbf{Global Circuit Architecture and Execution:}
The boolean gates are organised into a deep, feed-forward network of sequential layers. A critical architectural distinction exists between the global circuit inputs and the local inputs required by each gate. The circuit receives a global input vector $\mathbf{x} \in [0, 1]^{N_{in}}$, while individual gates operate with a fixed arity $k$ (fan-in). Signals propagate through the network via "wires," which, although they represent physical connections, are implemented as integer indices selecting specific bits from the previous layer. For example, consider a first layer of two arity-2 gates ($G_1, G_2$) processing a 4-bit global input $\mathbf{x} = [x_0, \dots, x_3]$. The wiring mechanism might assign indices $(0, 1)$ to $G_1$ and $(2, 3)$ to $G_2$. Consequently, $G_1$ computes $\text{LUT}_1(x_0, x_1)$ while $G_2$ independently computes $\text{LUT}_2(x_2, x_3)$. This indexing allows for arbitrary connectivity and signal duplication. The collective outputs of these gates then form a vector, which serves as the input for the subsequent layer.

\subsection{Graph Representation of Circuits}
To apply self-organising principles, we lift the circuit instance into a graph representation suitable for processing by a Graph Neural Network.
Given a circuit with a set of gates $G$ and wires $W$, we construct a graph $\mathcal{G} = (\mathcal{V}, \mathcal{E})$ where each node $v_i \in \mathcal{V}$ corresponds to a physical gate or an input/output node.

\textbf{Node State:}
The state of each node $v_i$ is a vector concatenation of functional and auxiliary features:
$$\mathbf{s}_i = [\boldsymbol{\ell}_i, \mathbf{m}_i, \mathbf{p}_i]$$
where:
\begin{itemize}
    \item $\boldsymbol{\ell}_i \in \mathbb{R}^{2^N}$ is the functional LUT logit vector defining the gate's logic.
    \item $\mathbf{m}_i \in \mathbb{R}^{d_{hidden}}$ is a latent memory vector allowing the meta-learner to store recurrent state across timesteps (hidden states).
    \item $\mathbf{p}_i$ is a sinusoidal positional encoding of the gate's normalised depth fraction $\frac{d_i}{D}$, where $d_i$ is the layer index and $D$ the total number of layers. By encoding depth as a ratio rather than an absolute index, this representation remains invariant to circuit scale: a gate at the midpoint of a 4-layer circuit receives the same encoding as one at the midpoint of a 10-layer circuit.
\end{itemize}
Optionally, a scalar local feedback signal $r_i$ (described in \S\ref{sec:node_loss}) can be appended to provide the meta-learner with per-node error context, yielding $\mathbf{s}_i = [\boldsymbol{\ell}_i, \mathbf{m}_i, \mathbf{p}_i, r_i]$.

\textbf{Connectivity:}
Crucially, the graph topology is dictated by the circuit wiring: 2 nodes are connected only if a wire exists between the corresponding gates. We enforce bidirectional message passing: if Gate A feeds Gate B, the graph contains edges $A \rightarrow B$ but also $B \rightarrow A$.
This allows gradients and error information to flow "backwards" (upstream) through the circuit topology during the self-organisation process, even though the logic signals only flow forward.

% \begin{figure}[h!]
%     \centering
%     \includegraphics[width=0.9\linewidth]{Figures/chap5/channels.png}
%     \caption{Schematic of a circuit, and of a specific node's state}
%     \label{ch5:fig:node_state}
% \end{figure}

\begin{figure}[ht]
    \centering
    % Input the tikz file relative to your main.tex
    % Use \resizebox if you need to scale the whole figure down to fit margins
    \resizebox{\linewidth}{!}{%
        \input{Figures/chap5/circuits.tikz}
    }
    \caption[Circuit Architecture and Graph Representation]{\textbf{Circuit Architecture and Graph Representation.} \textbf{(A)} The physical topology of the Differentiable Logic Gate Network. Gates (nodes) are arranged in layers and connected by wires (indices). \textbf{(B)} The graph representation of a single node state $\mathbf{h}_i$. This state vector concatenates the functional parameters (Logits) with the meta-learning variables (Memory, Position, Feedback) that the Graph Transformer utilizes to optimize the circuit.}
    \label{ch5:fig:circuit_schematic}
\end{figure}

% \subsection{Meta-Learning Architecture: Topology-Masked Transformer}
% \label{sec:architecture}
% To parameterise the update rule, we employ a single-block Transformer whose attention is restricted by the circuit topology. Each gate corresponds to a token, and the attention mask enforces that information flows only along the circuit's wires---combining the expressivity of self-attention with the locality of graph message passing. The model operates on the graph $\mathcal{G}$ defined in \cref{sec:LGNs}, processing all $N$ nodes in parallel to compute parameter updates.

% \textbf{Feature Projection:}
% The input node states $\{\mathbf{s}_i\}_{i \le N}$ (concatenating logits $\boldsymbol{\ell}_i$, hidden memory $\mathbf{m}_i$, positional encoding $\mathbf{p}_i$, and optionally $r_i$) are first normalised via LayerNorm \citep{ba2016layernormalization} to harmonise the heterogeneous feature scales, then projected into a high-dimensional latent space:
% \begin{equation}
%     \mathbf{z}^{in} = W_{in}\, \text{LN}(\mathbf{s})
% \end{equation}

% \textbf{Masked Self-Attention Block:}
% The core processing consists of a single Transformer block designed to respect the circuit topology. We enforce connectivity via a binary mask $M$, where $M_{ij} = 1$ if there is a wire between gates $j$ and $i$ (or if $i=j$), and $0$ otherwise. We adopt the \textbf{Pre-LN} Transformer convention \citep{xiong2020layernormalizationtransformerarchitecture}, applying LayerNorm \textit{before} each sub-layer rather than after. This is complemented by \textbf{QK-normalisation} \citep{dehghani2023scalingvisiontransformers22}, which normalises the query and key vectors \textit{after} projection to stabilise attention logits---particularly important in our setting where the same block is applied recurrently for many steps. To ensure stable training dynamics, we further employ \textbf{ReZero} regularisation \citep{bachlechner_rezero_2020}, where residual branches are gated by learnable scalars ($\alpha$) initialised to zero. This allows the block to initially act as an identity function, gradually learning to modify nodes. Altogether, the block computes an incremental refinement of the latents $\mathbf{z}^{in}$ :

% \begin{align}
%     \mathbf{a} &= \text{Attn}\!\left(\text{LN}_Q(\mathbf{z}^{in}),\; \text{LN}_{KV}(\mathbf{z}^{in}),\; \text{Mask}{=}M\right) && \text{(Pre-LN + Masked Attention)} \\
%     \mathbf{z}' &= \mathbf{z}^{in} + \alpha_{attn} \cdot \mathbf{a} && \text{(ReZero Residual)} \\
%     \mathbf{f} &= \text{MLP}\!\left(\text{LN}(\mathbf{z}')\right) && \text{(Pre-LN + FFN)} \\
%     \mathbf{z}^{out} &= \mathbf{z}' + \alpha_{ffn} \cdot \mathbf{f} && \text{(ReZero Residual)}
% \end{align}
% where $\text{LN}_Q$ and $\text{LN}_{KV}$ are separate LayerNorms for query and key/value inputs respectively, and MLP is a two-layer network with GELU activation.

% \textbf{Updates:}
% % Finally, the processed latent state $\mathbf{z}_i^{out}$ is decoded into specific parameter updates using two separate linear heads ($W_{l}, W_{m}$). These projections are also governed by ReZero gating to ensure the update magnitude is controllable and initially zero:
% A final LayerNorm is applied to the processed latent $\mathbf{z}_i^{out}$ before decoding into parameter updates via two separate linear heads ($W_{\ell}, W_{m}$), each governed by ReZero gating:

% % \begin{align}
% %     \Delta \mathbf{l}_i &= \alpha_{l} \cdot W_{l} \mathbf{z}_i^{out} && \text{(Logits Updates)} \\
% %     \Delta \mathbf{m}_i &= \alpha_{m} \cdot W_{m} \mathbf{z}_i^{out} && \text{(Latents Updates)}
% % \end{align}
% % These delta updates are added back to the current node state, effectively integrating the updates proposed by the meta-learner.

% \begin{align}
%     \hat{\mathbf{z}} &= \text{LN}_{out}(\mathbf{z}^{out}) \\
%     \Delta \boldsymbol{\ell} &= \alpha_{\ell} \cdot W_{\ell}\, \hat{\mathbf{z}} && \text{(Logit Updates)} \\
%     \Delta \mathbf{m} &= \alpha_{m} \cdot W_{m}\, \hat{\mathbf{z}} && \text{(Hidden State Updates)}
% \end{align}
% These delta updates are added back to the current node state, $\boldsymbol{\ell}_i \leftarrow \boldsymbol{\ell}_i + \Delta\boldsymbol{\ell}_i$ and $\mathbf{m}_i \leftarrow \mathbf{m}_i + \Delta\mathbf{m}_i$, while positional encodings remain fixed.


\subsection{Meta-Learning Architecture: Topology-Masked Transformer}
\label{sec:architecture}

We parameterise the update rule as a single Transformer block operating on $N$ gate tokens in parallel, with the circuit's connectivity enforced through an attention mask.

\textbf{Feature Projection.}
The node states $\mathbf{S} \in \mathbb{R}^{N \times d_{in}}$ (concatenating logits $\boldsymbol{\ell}_i$, memory $\mathbf{m}_i$, and position $\mathbf{p}_i$) are normalised and projected into a latent space:
\begin{equation}
    \mathbf{Z}^{(0)} = \text{LN}_{in}(\mathbf{S})\, W_{in}^T
\end{equation}

\textbf{Masked Attention Block.}
The latents are refined by a single Transformer block whose attention is restricted to wired neighbours via a binary mask $M \in \{0, -\infty\}^{N \times N}$.
We adopt \textbf{Pre-LN} normalisation \citep{xiong2020layernormalizationtransformerarchitecture} with \textbf{QK-normalisation} \citep{dehghani2023scalingvisiontransformers22} to stabilise attention logits across recurrent steps, and \textbf{ReZero} gating \citep{bachlechner_rezero_2020} ($\alpha$ scalars initialised near zero) so that the block initially acts as an identity:
\begin{align}
    \mathbf{Z}' &= \mathbf{Z}^{(0)} + \alpha_{\text{attn}} \cdot \text{MSA}\!\left(\text{LN}_Q(\mathbf{Z}^{(0)}),\; \text{LN}_{KV}(\mathbf{Z}^{(0)}),\; M\right) \\
    \mathbf{Z}'' &= \mathbf{Z}' + \alpha_{\text{ffn}} \cdot \text{MLP}\!\left(\text{LN}(\mathbf{Z}')\right) \\
    \mathbf{Z}_{out} &= \text{LN}_{final}(\mathbf{Z}'')
\end{align}
Separate LayerNorms for queries and keys/values decouple their statistics; the position-wise MLP ensures all cross-node communication is confined to the attention step.

\textbf{Output Updates.}
The output latents are decoded into parameter deltas via linear heads, again ReZero-gated:
\begin{align}
    \Delta \boldsymbol{\ell}_i &= \alpha_{\ell} \cdot \mathbf{Z}_{out,i}\, W_{\ell}^T, &
    \Delta \mathbf{m}_i &= \alpha_{m} \cdot \mathbf{Z}_{out,i}\, W_{m}^T
\end{align}
applied as residual updates: $\boldsymbol{\ell}_i \leftarrow \boldsymbol{\ell}_i + \Delta\boldsymbol{\ell}_i$, $\mathbf{m}_i \leftarrow \mathbf{m}_i + \Delta\mathbf{m}_i$.

\textbf{Per-Node Error Feedback.}
\label{sec:node_loss}
As described, the architecture has no direct information about how well the circuit performs on its task. For fixed-wiring experiments this is sufficient; however, when generalising across \emph{random wirings}, we found it critically insufficient. We augment each output gate's state with a scalar feedback signal $r_i$---the average absolute residual across a batch of task evaluations:
\begin{equation}
    r_i = \frac{1}{|\mathcal{D}|}\sum_{(\mathbf{x}, \mathbf{y}) \in \mathcal{D}} |y_i - \hat{y}_i(\mathbf{x})|
\end{equation}
where $\hat{y}_i$ is the soft output of gate $i$ on input $\mathbf{x}$. This signal is recomputed at every recurrent step; non-output gates receive $r_i = 0$. Empirically, this simple mechanism is essential for wiring-agnostic optimisation: the signal propagates upstream through recurrent attention steps, allowing interior gates to adjust in response to downstream error---a decentralised form of credit assignment. We note that this still leaves the meta-learner \textit{blind} to the actual input--output data; cross-attention to task tokens (in the spirit of Perceiver IO \citep{jaegle2022perceiveriogeneralarchitecture}) is a promising avenue for removing this limitation, and is the subject of ongoing work.

\textbf{Recurrence as Depth.}
The architecture employs a \textit{single} attention block applied recurrently for $T$ steps:
\begin{equation}
    \mathbf{s}_i^{(t+1)} = \mathbf{s}_i^{(t)} + \Delta\mathbf{s}_i^{(t)}, \quad t = 1, \dots, T
\end{equation}
where $\Delta\mathbf{s}_i^{(t)}$ is computed by the same shared-weight block at every step---functionally equivalent to a $T$-layer weight-tied residual network. Expressivity comes through \textit{iterated refinement}: each step nudges each gate based on its current neighbourhood, and the accumulation of many such nudges achieves globally coordinated solutions. The topology mask imposes a strict ``speed of light'': information propagates at most one hop per step, so a gate $k$ edges away requires $\geq k$ steps to exert influence. Global coordination thus emerges from iterated local interactions---a deliberate inductive bias faithful to the NCA paradigm, ensuring the policy is fundamentally \textit{local} and that global solutions arise through repeated application of local rules over time.

\textbf{Scale-Free Architecture.}
Because the Transformer operates at the node level with shared weights, its parameters are independent of circuit size---the same update rule applies to 20 or 200 gates, 3 or 10 layers. Only the topology mask $M$ must be recomputed from the wiring. Combined with normalised positional encodings, this endows the architecture with \textit{scale-freedom}: a trained policy can, in principle, be deployed on circuits of different size without retraining.



\subsection{Training Framework: Pool-Based Meta-Learning}
We treat the circuit optimisation problem as a dynamic system trained via Backpropagation Through Time (BPTT).

\textbf{Pool Mechanism:}
As in the original NCA paper, we maintain a persistent \textbf{Graph Pool} $\mathcal{P} = \{ \mathcal{G}_1, \dots, \mathcal{G}_K \}$ containing $K$ circuit instances being continually optimized. By persisting state across many training steps, this training mechanism favours policies that achieve and maintain homeostasis, as circuits must remain stable over potentially long durations. To ensure diversity and prevent stagnation, a portion of the pool is periodically reset, injecting ``fresh'' unoptimized circuits back into the process. The reset interval is computed automatically based on other parametrs, namely the batch size and number of message steps applied per batch, so that circuits have been optimized for 128 steps on average before reset. 
% While this can seem low to reach complete homeostasis we find that  

\textbf{Inner Loop (Inference --- ``Growth''):}
At each training step, a batch of circuits is sampled from the pool. A batch of Boolean input--output pairs $(\mathbf{X}, \mathbf{Y})$ is drawn for the current task. The Transformer is then applied iteratively for $T$ steps via a \texttt{scan} operation. Crucially, the circuit is functionally executed \textit{at every step}: at each tick $t$, the updated LUT logits are extracted, the circuit is run on $(\mathbf{X}, \mathbf{Y})$, and the resulting per-node residuals $r_i$ are written back into the graph before the next application. This creates a closed feedback loop where the meta-learner observes the consequences of its updates in real time. 
Regarding the temporal horizon, while stochastic update depths are possible (variable $T$ per sample of the batch), we found that backpropagating through a fixed number of steps yielded the best results. The optimal truncation horizon ($T$) varied by task complexity: simple experiments (e.g., fixed wiring) converged with short horizons ($T=5$), whereas complex, wiring-agnostic optimization required significantly deeper unrolling ($T=50$).

\textbf{Outer Loop (Optimisation):}
After the $T$-step scan, the circuit loss $\mathcal{L}$ (Binary Cross Entropy between circuit output and target $\mathbf{Y}$) at a selected step $t^{*}$ is used as the training objective. We support both fixed evaluation at $t^{*} = T$ and stochastic evaluation where $t^{*}$ is drawn uniformly from $\{T_{\min}, \dots, T\}$ per sample---the latter encouraging the policy to produce functional circuits at any point during the optimisation, not only at the final step. Gradients of $\mathcal{L}$ with respect to the NCA parameters are computed via BPTT through the differentiable scan and circuit execution, and parameters are updated using Adam \citep{kingma2014adam}.

\begin{figure}
    \centering
    \includegraphics[width=0.95\linewidth]{Figures/chap5/minimalist.pdf}
    \caption[Summary of training procedure]{\textbf{Summary of meta-learning training framework.} Step 1: a pool of on-going optimization circuit graphs is sampled to get a (meta) batch. Step 2: the sampled graphs are optimized for a fixed horizon $T$ using the NCA message passing scheme. Step 3: the logits are extracted from the graphs, and the resulting circuits are functionally executed on the task at hand. The resulting loss $\mathcal{L}$ is used to backpropagate gradients and update the NCA policy}
    \label{ch5:fig:training}
\end{figure}


\section{Experiments}
\label{sec:experiments}

To systematically evaluate the capabilities of the self-organising circuit framework, we designed a curriculum of experiments acting as a stress-test for the meta-learner. We progress from idealised conditions—growing a circuit on a known, fixed architecture—to increasingly hostile environments characterised by physical damage, and finally to creating generalist policies capable of navigating unseen, random topologies.

\subsection{Regime I: \textit{Ex Nihilo} Growth on Fixed Topologies}
\label{subsec:fixed_growth}

Our initial experiment serves as a sanity check: can the NCA policy replace standard backpropagation to configure a functional circuit from scratch?

\textbf{Setup:}
In this regime, the Graph Pool is initialised with circuits sharing a single, fixed wiring architecture (Fixed Topology). The circuits are initialised with random, non-functional LUT configurations. The meta-optimiser is tasked with "growing" functionality—guiding the system from this random state to a solution that satisfies a specific Boolean task (e.g., binary multiplication).

\textbf{Pool Dynamics:}
Following the standard NCA protocol, we utilise a persistent pool to stabilise training. However, unlike the original NCA implementation where the pool is introduced late to ensure homeostasis, we employ it from the outset. It is important to note that in this fixed-topology regime, the meta-learner is effectively overfitting to a specific graph structure; the "solution" is a policy specialised for this specific wiring diagram.

\textbf{Objective:}
The primary goal here is to demonstrate convergence. We validate that the local update rules can successfully navigate the Boolean search space to find functional configurations, achieving accuracy comparable to standard global backpropagation.

\subsection{Regime II: Robustness and Adaptive Resilience}
\label{subsec:robustness}

Having established that the system can construct functional circuits, we next investigate its ability to maintain them in a hostile environment. We introduce structural perturbations into the training loop, testing the hypothesis that exposing the meta-learner to damage will promote the emergence of adaptive resilience.

\textbf{Damage Modelling:}
To simulate hardware faults, we introduce two distinct categories of damage during the meta-learning process:

\begin{enumerate}
    \item \textbf{Recoverable Damage (Soft Errors):} To simulate temporary upsets (e.g., bit flips), specific gates are forced to a faulty value (clamped to -10 logits) but remain active in the computational graph. The optimiser can potentially recover the correct logic for these gates in subsequent timesteps.
    \item \textbf{Permanent Damage (Stuck-at Faults):} To simulate catastrophic hardware failure, we apply permanent gate knockouts. Damaged nodes are clamped to a faulty value and permanently banned from receiving updates. Crucially, they remain in the topology and continue to emit messages, signalling their faulty status to neighbours. This forces the system to recognise dead components and "work around" them, re-routing information through the remaining functional substrate.
\end{enumerate}

\textbf{Training Protocol:}
We incorporate these failure modes directly into the training process. At each step, a probability of failure is calculated such that approximately 10\% of the gates fail by the time the circuit is reset. This creates a "survival of the fittest" pressure: the policy must not only solve the task but doing so while anticipating the loss of resources.

\textbf{OOD Evaluation (The "Shotgun" Test):}
To test robustness Out-Of-Distribution (OOD), we subject the trained circuits to a "shotgun" damage pattern during evaluation. Instead of the gradual failure seen during training, 5\% of the gates are instantly destroyed in a single step, repeated twice over the evaluation window.
We hypothesise that training with damage acts as a rigorous form of regularisation, preventing brittle solutions and encouraging the system to develop both redundancy (minimising the immediate drop in performance) and adaptation (recovering loss after the damage event).

\subsection{Regime III: Pure Maintenance and Homeostasis}
\label{subsec:maintenance}

In the previous regimes, the NCA acted as both architect (growing the circuit) and repairman. Here, we isolate the "repairman" role to evaluate the system's capacity for pure maintenance.

\textbf{Setup:}
Instead of seeding the pool with random initialisations, we seed it with \textit{functional} circuits that have been pre-optimised using standard backpropagation. The NCA is not tasked with finding the solution, but solely with preserving high-performance functionality in the face of the damage models described in \cref{subsec:robustness}. This experiment parallels the biological reality where maintenance mechanisms operate on structures already formed by development.

\subsection{Regime IV: Generalisation to Random Topologies}
\label{subsec:random_wiring}

The most challenging setting removes the architectural constraints entirely, denying the meta-learner the ability to memorise a specific graph structure.

\textbf{Setup:}
We transition to a \textbf{Random Topology} regime, where every circuit in the pool possesses a unique, randomly generated wiring diagram.
Consequently, the meta-learner cannot overfit to "Gate A connected to Gate B." Instead, it must learn a truly topological, wiring-agnostic policy that generalises to unseen architectures.

\textbf{Implications:}
We anticipate this regime to be significantly harder than Fixed Topology training. Success in this regime would imply the discovery of a "universal" local update rule for Boolean logic—a policy capable of entering an arbitrary, random arrangement of gates and organising them into a functional computer. While we expect the performance to lag behind global backpropagation, the ability to make progress in this regime represents a step toward truly general-purpose self-organising matter.

\subsection{Regime V: Scale-Free Optimisation}
\label{subsec:scalability}

Finally, we leverage the decentralised nature of our architecture to explore scale capabilities. Because the Graph Transformer shares weights across all nodes and operates on local neighbourhoods, the learned update rule is theoretically independent of the circuit size. 


\section{Discussion}

\subsection*{Summary}

We have presented a framework for self-organising digital circuits that replaces global backpropagation with a decentralised, learned update rule. By extending the NCA paradigm from grid-based pattern formation to functional logic generation on arbitrary graphs, we demonstrated that a single topology-masked Transformer block, applied recurrently, can learn to configure Boolean circuits to satisfy computational tasks.

On fixed topologies, the meta-learner assembles functional circuits from scratch with accuracy comparable to standard backpropagation, and maintains homeostasis over extended time horizons. When stochastic gate failures are introduced during training, the policy develops emergent robustness---exhibiting both redundancy (reduced performance drops) and active adaptation (recovery after damage)---that generalises to unseen, out-of-distribution fault patterns. Notably, repair does not consist of returning to a memorised pre-damage configuration; instead, the system navigates the solution manifold to discover \textit{new} functional arrangements, consistent with the biological paradigm of adaptive plasticity.

Extending to random topologies---where every circuit in the pool has a unique wiring diagram---we found that providing per-node error feedback (\S\ref{sec:node_loss}) is essential for the meta-learner to distinguish between structurally distinct circuits. While this regime remains significantly more challenging than fixed wiring, the policy makes meaningful progress toward topology-agnostic optimisation. 
% The scale-free nature of the architecture further allows a policy trained on circuits of a given size to be deployed on larger or differently structured circuits without retraining.

\subsection*{Limitations}

Several limitations of the current framework should be acknowledged. First, the recurrent unrolling required for complex tasks (up to $T=25$ steps for random wiring) makes training expensive: backpropagation through time across many differentiable circuit executions demands substantial memory and compute, despite gradient checkpointing. Second, the per-node error feedback $r_i$, while effective, is a coarse signal---it provides only a scalar average residual per output gate, leaving the meta-learner without direct access to the input--output data that generated it. This limits the policy's ability to perform true in-context reasoning about the task structure. Third, the current positional encoding scheme (normalised depth fraction) captures only vertical position in the circuit DAG, offering no information about the local graph topology or the structural role of each gate.

\subsection*{Future Directions}

\textbf{Enriching the Graph Representation.}
Our current architecture is deliberately minimal: the circuit topology enters only through the attention mask. While this simplicity ensures generality, it also means the meta-learner must reconstruct structural information about its local neighbourhood entirely through recurrent propagation---a slow process when the graph is large. Techniques from the Graph Transformer literature offer principled ways to accelerate this, while preserving the scale-free property that is central to our design.

\textit{Random Walk Structural Encoding} (RWSE) \citep{dwivedi2021generalizationtransformernetworksgraphs} is a particularly natural fit: for each node, one computes the diagonal entries of the $k$-step random walk matrix $\mathbf{A}^k$ for $k = 1, \dots, K$, yielding a $K$-dimensional structural fingerprint that captures local topology (e.g., whether a gate sits in a tightly clustered region, at a bottleneck, or at the end of a long chain). Crucially, RWSE is permutation-equivariant, inexpensive to compute from the adjacency matrix, and---because it characterises \textit{local structure} rather than absolute position---transfers naturally across circuits of different sizes. Similarly, encoding each gate's fan-out degree (the number of downstream consumers) would provide an immediate, zero-cost signal about structural importance.

Another promising direction is to inject \textit{edge-type information} into the attention computation. Currently, the model cannot distinguish between forward wires (carrying logic signals downstream) and backward edges (carrying error information upstream). Introducing learned scalar biases per edge type---forward, backward, and self-loop---directly into the attention logits would allow the meta-learner to treat upstream and downstream messages differently, at the cost of only a handful of additional parameters.

All of these enrichments respect the key constraint of scale-freedom: they are defined locally, do not depend on absolute graph size, and can be recomputed for any circuit at inference time.

\textbf{Toward True In-Context Learning.}
A more fundamental extension concerns the meta-learner's access to task information. As discussed in \S\ref{sec:node_loss}, the current system receives only a scalar residual per output gate---a compressed shadow of the full input--output relationship that generated it. This prevents the emergence of genuine in-context learning, where the optimiser would adapt its strategy based on the specific structure of the data.

A natural remedy is to augment the architecture with \textit{cross-attention} to the task data itself, in the spirit of Perceiver IO \citep{jaegle2022perceiveriogeneralarchitecture}. In this setting, the input--output pairs $(\mathbf{x}, \mathbf{y})$ would be tokenised and made available as a separate key-value set that the gate tokens can attend to at each recurrent step. This would allow the meta-learner to directly inspect which input patterns cause which output errors, enabling it to make targeted, data-informed corrections rather than relying solely on aggregate residual signals. Stochastic subsampling of data tokens---attending to a random subset at each step---could further promote incremental in-context learning while controlling the computational cost of cross-attention. Preliminary experiments with this Perceiver-style architecture are ongoing, and early results on simple tasks are encouraging.


\textbf{Sparse Attention as a Computational Advantage.}
A practical limitation of the current implementation is that the topology mask is applied as a dense $N \times N$ matrix: the full attention score matrix is computed, and only then are disallowed entries zeroed out. This means the computational cost remains $\mathcal{O}(N^2 d)$ per step, even though the number of active (non-masked) entries scales as $\mathcal{O}(E)$, the number of edges in the circuit graph. For bounded-arity circuits, $E$ grows linearly with $N$, so the gap between the dense cost and the true information-theoretic cost is quadratic.

To make this concrete: a typical circuit in our experiments has $N=264$ nodes (12 inputs, three hidden layers of 96, 96, and 48 gates with arity $k=4$, and 12 outputs) and $E=2{,}016$ bidirectional edges (including self-loops). The dense attention matrix contains $264^2 = 69{,}696$ entries, of which only $2{,}016$ are non-zero---a density of \textbf{2.9\%}. Over 97\% of the attention computation is wasted on entries that are guaranteed to be masked out.

This inefficiency is compounded by the recurrent architecture. Each of the $T$ optimisation steps recomputes the full dense attention, and backpropagation through time requires storing these matrices at every step. For our setup with $T=50$ steps, the dense backward pass stores $50 \times 69{,}696 \approx 3.5$M attention entries; with true sparse attention this would be $50 \times 2{,}016 \approx 100$K---a \textbf{$\mathbf{35\times}$ memory reduction}. For scale-free deployment to larger circuits, the savings grow quadratically: doubling the circuit size quadruples the dense cost but only doubles the sparse cost.

Importantly, the sparsity here is not approximate: it is \textit{exact and known a priori} from the circuit wiring. This makes it fundamentally different from the soft or learned sparsity pursued in efficient Transformer research \citep{tay2022efficienttransformerssurvey}. The topology mask defines a precise, fixed communication graph, and no information is lost by restricting attention to it. Implementing this as a native sparse operation---computing attention weights only along existing edges, in the style of Graph Attention Networks \citep{velickovic_graph_2018} but within the Transformer framework---would reduce both compute and memory to scale with the circuit's \textit{wiring} rather than its \textit{size}. This would make it practical to apply the learned policy to circuits with thousands of gates, where dense attention becomes prohibitive but the per-node edge count remains bounded by the arity.

Finally, there is a suggestive connection to hardware deployment. On reconfigurable substrates such as FPGAs, sparse communication patterns can be physically routed into the fabric, instantiating only the compute paths that the topology requires. The same structural sparsity that currently goes unexploited on GPUs could become a natural efficiency advantage on the target hardware, aligning the computational structure of the meta-learner with the physical structure of the circuit it optimises.

Together, these directions---richer graph encodings, direct data access, and sparse-native computation---point toward a meta-learner that is not only topology-agnostic and scale-free, but also \textit{task-aware} and \textit{computationally efficient}: a truly general-purpose, decentralised optimiser for digital substrates.